% p072_1; source PDF page 72; manually verified editable table.
\begin{tikzpicture}[x=1cm,y=-1cm,inner sep=0pt,outer sep=0pt]
\path[use as bounding box] (0,0) rectangle (16.000,5.985);
\draw[line width=.4pt] (0.000,0.000) rectangle (4.600,0.855);
\node[align=center,text width=4.340cm,font={\linespread{1}\fontsize{10}{13.50}\selectfont\bfseries}] at (2.300,0.427) {罪过形式};
\draw[line width=.4pt] (4.600,0.000) rectangle (10.400,0.855);
\node[align=center,text width=5.540cm,font={\linespread{1}\fontsize{10}{13.50}\selectfont\bfseries}] at (7.500,0.427) {认识因素};
\draw[line width=.4pt] (10.400,0.000) rectangle (16.000,0.855);
\node[align=center,text width=5.340cm,font={\linespread{1}\fontsize{10}{13.50}\selectfont\bfseries}] at (13.200,0.427) {意志因素};
\draw[line width=.4pt] (0.000,0.855) rectangle (4.600,1.710);
\node[align=left,text width=4.340cm,font={\linespread{1}\fontsize{10}{13.50}\selectfont}] at (2.300,1.282) {直接故意};
\draw[line width=.4pt] (4.600,0.855) rectangle (10.400,1.710);
\node[align=left,text width=5.540cm,font={\linespread{1}\fontsize{10}{13.50}\selectfont}] at (7.500,1.282) {认识到必然或可能发生};
\draw[line width=.4pt] (10.400,0.855) rectangle (16.000,1.710);
\node[align=left,text width=5.340cm,font={\linespread{1}\fontsize{10}{13.50}\selectfont}] at (13.200,1.282) {希望（赞成票）};
\draw[line width=.4pt] (0.000,1.710) rectangle (4.600,2.565);
\node[align=left,text width=4.340cm,font={\linespread{1}\fontsize{10}{13.50}\selectfont}] at (2.300,2.138) {间接故意};
\draw[line width=.4pt] (4.600,1.710) rectangle (10.400,2.565);
\node[align=left,text width=5.540cm,font={\linespread{1}\fontsize{10}{13.50}\selectfont}] at (7.500,2.138) {认识到可能发生};
\draw[line width=.4pt] (10.400,1.710) rectangle (16.000,2.565);
\node[align=left,text width=5.340cm,font={\linespread{1}\fontsize{10}{13.50}\selectfont}] at (13.200,2.138) {放任（弃权票）};
\draw[line width=.4pt] (0.000,2.565) rectangle (4.600,3.420);
\node[align=left,text width=4.340cm,font={\linespread{1}\fontsize{10}{13.50}\selectfont}] at (2.300,2.992) {过于自信的过失};
\draw[line width=.4pt] (4.600,2.565) rectangle (10.400,3.420);
\node[align=left,text width=5.540cm,font={\linespread{1}\fontsize{10}{13.50}\selectfont}] at (7.500,2.992) {认识到可能发生，本应避免};
\draw[line width=.4pt] (10.400,2.565) rectangle (16.000,3.420);
\node[align=left,text width=5.340cm,font={\linespread{1}\fontsize{10}{13.50}\selectfont}] at (13.200,2.992) {不想发生（反对票）};
\draw[line width=.4pt] (0.000,3.420) rectangle (4.600,4.275);
\node[align=left,text width=4.340cm,font={\linespread{1}\fontsize{10}{13.50}\selectfont}] at (2.300,3.848) {疏忽大意的过失};
\draw[line width=.4pt] (4.600,3.420) rectangle (10.400,4.275);
\node[align=left,text width=5.540cm,font={\linespread{1}\fontsize{10}{13.50}\selectfont}] at (7.500,3.848) {没有预见到，但应当预见到};
\draw[line width=.4pt] (10.400,3.420) rectangle (16.000,4.275);
\node[align=left,text width=5.340cm,font={\linespread{1}\fontsize{10}{13.50}\selectfont}] at (13.200,3.848) {不想发生};
\draw[line width=.4pt] (0.000,4.275) rectangle (4.600,5.130);
\node[align=left,text width=4.340cm,font={\linespread{1}\fontsize{10}{13.50}\selectfont}] at (2.300,4.703) {意外事件};
\draw[line width=.4pt] (4.600,4.275) rectangle (10.400,5.130);
\node[align=left,text width=5.540cm,font={\linespread{1}\fontsize{10}{13.50}\selectfont}] at (7.500,4.703) {没有预见到，且无法预见到};
\draw[line width=.4pt] (10.400,4.275) rectangle (16.000,5.130);
\node[align=left,text width=5.340cm,font={\linespread{1}\fontsize{10}{13.50}\selectfont}] at (13.200,4.703) {不想发生};
\draw[line width=.4pt] (0.000,5.130) rectangle (4.600,5.985);
\node[align=left,text width=4.340cm,font={\linespread{1}\fontsize{10}{13.50}\selectfont}] at (2.300,5.558) {不可抗力};
\draw[line width=.4pt] (4.600,5.130) rectangle (10.400,5.985);
\node[align=left,text width=5.540cm,font={\linespread{1}\fontsize{10}{13.50}\selectfont}] at (7.500,5.558) {预见到，但无法避免};
\draw[line width=.4pt] (10.400,5.130) rectangle (16.000,5.985);
\node[align=left,text width=5.340cm,font={\linespread{1}\fontsize{10}{13.50}\selectfont}] at (13.200,5.558) {不想发生};
\end{tikzpicture}
